\subsection{The spectrum and the first functions eliminated by data}

\begin{proposition}[Paired spectrum and minimum-energy signs]
\label{prop:paired-spectrum}
Let $\lambda_r$ be the eigenvalue of $M^{-1}K$ on degree-$r$ parities.
For the kernel~\eqref{eq:arcsine-kernel},
\[
 \lambda_0=\lambda_1>\lambda_2=\lambda_3>
 \lambda_4=\lambda_5>\cdots>0,
\]
with a final unpaired $\lambda_d$ when $d$ is even. The minimum of
$E_K(g)$ over all Boolean functions is $a^2/(2\lambda_0)$, attained
exactly by the constants and signed coordinate functions
$g\in\{1,-1,x_1,-x_1,\ldots,x_d,-x_d\}$.
For every finite width, all these functions have identical
unconditioned full-margin prior probabilities.
\end{proposition}

\begin{proof}
Put $D=d+1$ and write $F$ for the $(L-1)$-fold composition of
$t\mapsto2\arcsin(t)/\pi$. Its power series is
$F(t)=\sum_{m\text{ odd}}c_mt^m$, with every $c_m>0$ and
$\sum_m c_m=1$. Positivity follows from the arcsine series and is
preserved under composition. First consider the kernel $F(u\cdot v/D)$
on the full $D$-dimensional cube. The degree-$r$ Fourier eigenvalue
of $(u\cdot v/D)^m$ is the probability $p_m(r)$ that a simple random
walk, starting at one cube vertex and flipping a uniformly chosen
coordinate at each step, reaches a specified vertex at Hamming
distance $r$ after $m$ steps. This follows by expanding the power:
a term survives Fourier averaging exactly when the coordinates
chosen an odd number of times form the specified support.

These probabilities decrease with distance within each parity class.
Indeed, for $0\le r\le D-2$,
\begin{align*}
 D\bigl[p_{m+1}(r)-p_{m+1}(r+2)\bigr]
 ={}&r\bigl[p_m(r-1)-p_m(r+1)\bigr]\\
 &+(D-r-2)\bigl[p_m(r+1)-p_m(r+3)\bigr],
\end{align*}
where terms with zero coefficients at the boundaries are omitted.
Induction starts from the point mass at time zero. Since
$p_r(r)>0=p_r(r+2)$ and $c_r>0$ for odd $r$, the positive odd-degree
eigenvalues of the full-cube kernel decrease strictly; its even-degree
eigenvalues vanish. Restricting the distinguished coordinate to $+1$
makes the reduced degree-$r$ eigenvalue the sum of the full-cube
degree-$r$ and degree-$(r+1)$ eigenvalues. This gives the paired
ordering stated.

The largest matrix eigenvalue of $K$ is $M\lambda_0$. Every
full-margin output satisfies
\[
 \frac12z^TK^{-1}z\ge\frac{\|z\|_2^2}{2M\lambda_0}
                     \ge\frac{a^2}{2\lambda_0}.
\]
Equality requires both $|z(x)|=a$ everywhere and
$z/a$ in the span of $1,x_1,\ldots,x_d$. An affine function
$c+\sum_j b_jx_j$ is Boolean only if it is a constant or signed
coordinate: expanding its square shows $b_ib_j=0$ for $i\ne j$
and $cb_i=0$, so at most one coefficient is nonzero, with magnitude
one. These functions attain equality, completing the characterization.

To prove the finite-width equality of prior weights, exchange
columns $0$ and $j$ of the first-layer weight matrix. Its Gaussian
law is unchanged. Let $P_j$ denote this exchange and let $T_j$ be
the cube bijection defined by $(T_jx)_j=x_j$ and
$(T_jx)_i=x_jx_i$ for $i\ne j$. Then
$P_j(1,x)=x_j(1,T_jx)$. The network as a function of its augmented
input is odd, since all activations are odd and there are no later
biases. Thus the transformed output is $x_j z_N(T_jx)$.
The all-positive margin orthant is mapped bijectively to that of
$x_j$, proving equality. Flipping the readout weights proves
equality with the negative functions as well.
\end{proof}

\begin{corollary}[A posterior-error lower bound from surviving coordinates]
\label{cor:surviving-coordinates}
Fix $0<\gamma<1$ and let $f=\chi_U$ be a parity of size $k\ge2$.
If any constant or signed coordinate agrees with a training sample
$S$, the posterior converges to the uniform distribution over all
such compatible functions. In particular,
\[
 \lim_{N\to\infty}R_N(S,f)=\frac12.
\]
For $d>k$ and $n\ge1$ independent uniform training examples,
\begin{equation}\label{eq:coordinate-sample-bound}
 \liminf_{N\to\infty}\mathbb E_S R_N(S,f)
 \ge\frac12\left[1-\left(1-2^{1-n}\right)^{d-k}\right].
\end{equation}
\end{corollary}

\begin{proof}
If such a function survives, the compatible absolute minimizers in
Proposition~\ref{prop:paired-spectrum} exhaust the compatible
minimum-energy functions. Theorem~\ref{thm:kernel-posterior} forces
the posterior onto these functions, and their exact equal prior
weights make the conditional distribution uniform. Each has uniform error exactly
$1/2$ against a parity of degree at least two, independently of
the mixture weights. This proves the first statement.

Conditional on all samples in coordinates belonging to $U$, the
target labels are fixed. Each of the $d-k$ irrelevant coordinates
has an independent uniform sample column in $\{-1,1\}^n$. One of
its two signs agrees with all labels with probability $2^{1-n}$.
These events are conditionally independent. The probability that
at least one such coordinate survives is consequently
$1-(1-2^{1-n})^{d-k}$. On this event the limiting error is $1/2$;
on its complement error is nonnegative. Fatou's lemma proves the
bound.
\end{proof}

\begin{lemma}[Parity correlation with a halfspace]
\label{lem:halfspace-parity}
If $h$ is a Boolean halfspace on $\{-1,1\}^d$ and $|U|=k\ge1$, then
$|\mathbb E[\chi_U h]|\le1/\sqrt{k}$. Either inclusion or exclusion
of the threshold boundary is allowed.
\end{lemma}

\begin{proof}
Condition on coordinates outside $U$ and flip selected input
coordinates so that the halfspace is monotone in the remaining $k$
coordinates. These flips change the parity only by a constant sign.
For a monotone Boolean function $h$, its discrete derivative
$D_i h=(h(x_i=1)-h(x_i=-1))/2$ is nonnegative. Thus, for every
$i\in U$,
\[
 |\widehat h(U)|
 =|\mathbb E[\chi_{U\setminus\{i\}}D_ih]|
 \le\mathbb E D_i h=\widehat h(\{i\}).
\]
Summing over the $k$ coordinates and using Cauchy--Schwarz and
Parseval gives $k|\widehat h(U)|\le\sqrt{k}$.
Average the conditional estimate over the other coordinates.
\end{proof}

\begin{corollary}[Failure to attain posterior confidence $2/3$]
\label{cor:confidence-groundstates}
Under the hypotheses of Corollary~\ref{cor:surviving-coordinates},
if any constant or signed coordinate is consistent with $S$, then
\[
 \liminf_{N\to\infty}F_N(S,f)
 \ge b_k:=\max\left\{\frac14,\frac12-\frac1{2\sqrt{k}}\right\}.
\]
If there are $m$ such compatible functions, the additional bound
$\liminf_N F_N(S,f)\ge1-9/m$ holds.
Consequently, when $d>k$ and $n\ge1$,
\[
 \liminf_{N\to\infty}\mathbb E_S F_N(S,f)
 \ge b_k\left[1-\left(1-2^{1-n}\right)^{d-k}\right].
\]
\end{corollary}

\begin{proof}
The bound $1/4$ follows from
\eqref{eq:confidence-risk-comparison} and
Corollary~\ref{cor:surviving-coordinates}. To obtain the stronger
bound for large $k$, write $\epsilon_N\to0$ for the posterior mass
outside the consistent absolute minimizers. The mean Boolean
prediction conditioned on those minimizers is an affine function
$\widetilde m_N(x)=c_N+\sum_jb_{N,j}x_j$. The full posterior mean
$m_N(x)$ satisfies
$\|m_N-\widetilde m_N\|_\infty\le2\epsilon_N$.
Confidence success is the event $f(x)m_N(x)>1/3$, which implies
$f(x)\widetilde m_N(x)>t_N:=1/3-2\epsilon_N>0$ for large $N$.

For any affine $m$ and any $t>0$, set
$w(x)=\mathbf1\{m(x)>t\}-\mathbf1\{m(x)<-t\}$.
This is the average of two Boolean halfspaces, using the appropriate
boundary conventions. By Lemma~\ref{lem:halfspace-parity},
$|\mathbb E fw|\le1/\sqrt{k}$. Moreover
\[
 \Pr\{fm>t\}=\frac12\mathbb E(|w|+fw)
          \le\frac12+\frac1{2\sqrt{k}}.
\]
Apply this to $m=\widetilde m_N$ and $t=t_N$. This proves the
pointwise-in-sample bound without assuming convergence of the
weights among tied minimizers. The expected bound follows using
the event probability calculated in
Corollary~\ref{cor:surviving-coordinates} and Fatou's lemma.

For the additional bound, equal weights make
$\widetilde m_N=m^{-1}\sum_{j=1}^m g_j$ independent of $N$.
Distinct signed constants and coordinates are orthogonal unless
they are negatives of one another. Hence
$\mathbb E\widetilde m_N^2\le1/m$. Confidence success implies
$|\widetilde m_N|>t_N$, so Markov's inequality bounds its fraction
by $1/(mt_N^2)\to9/m$.
\end{proof}

More explicitly, put $D=d-k$ and $p=2^{1-n}$ for $n\ge1$.
The number of surviving irrelevant coordinates is
$\operatorname{Bin}(D,p)$. For any integer $t>9$, the same proof gives
\[
 \liminf_N\mathbb E_S F_N(S,f)
 \ge (1-9/t)\Pr\{\operatorname{Bin}(D,p)\ge t\}.
\]
Thus confidence failure can be close to one while many spurious
coordinates remain compatible with the sample.

These results establish a necessary sample scale of order
$\log(d-k)$ for posterior confidence in this regime. They are not
exponential-in-$k$ bounds, nor do they assert the same error for
posterior majority.
The spectrum also gives
$E_K(\chi_U)\ge(a^2/2)\binom dk$, since
$\sum_{r=0}^d\binom dr\lambda_r=1$. Even an exponentially large
parity energy when $d\ge2k$ is a prior statement rather than a
sample lower bound.
